\chapter{PROBLEMS ENCOUNTERED AND PROPOSED SOLUTIONS}
\thispagestyle{empty}

\section{Methodological Problems}

The most significant problem encountered during this phase was methodological.
The earlier evaluation branch used a directed outgoing-edge split while the
structural models, especially Node2Vec, ultimately operated on undirected graph
structure. This created a mismatch between the benchmark definition and the way
the structural signal was consumed by the model. The effect was especially
important for Node2Vec, whose earlier dominance became questionable once the
evaluation branch was redesigned.

A second methodological issue concerned fairness across model families. Reusing
old GCN and GraphSAGE caches alongside the revised Node2Vec branch would have
created an apples-to-oranges comparison. This was avoided by creating separate
retraining notebooks for the revised structural branch.

\section{Implementation and Engineering Problems}

Several practical problems also emerged during implementation:
\begin{itemize}
    \item one evaluation notebook became malformed at the JSON level even though
    the code logic itself was still meaningful,
    \item outdated caches in the old evaluation branch no longer matched the
    corrected logic,
    \item early stopping in the first draft of the revised GNN notebooks used
    the wrong edge set and therefore risked test leakage,
    \item the web application initially used only a subset of models and needed
    to be aligned with the full experimental family,
    \item title handling and search interaction in the interface required
    cleanup because Wikipedia article identifiers contain underscores while the
    user-facing interface should display natural titles,
    \item the experimental local LLM layer for learning-path generation proved
    much less stable than the ranked retrieval layer. Small-model outputs were
    often inconsistent, over-concentrated in a single category, or weakly
    justified at the semantic level.
\end{itemize}

\section{Adopted or Planned Solutions}

The solutions adopted in this phase are now part of the project baseline:
\begin{itemize}
    \item a new phase-0 notebook creates an undirected structural branch and
    exports leakage-free Node2Vec caches,
    \item new GCN and GraphSAGE notebooks retrain the structural GNN baselines
    on the same revised branch,
    \item the final ten-model evaluation is centralized in a dedicated notebook
    and cached under a separate namespace,
    \item stale caches from invalid branches were deleted or isolated instead of
    being silently reused,
    \item the web application was updated to expose all tested models and to
    provide a more stable search and preview experience.
\end{itemize}

The LLM-based learning-path layer also had to be redesigned during this phase.
The first organizer attempt asked the local model to produce a full structured
JSON path over the entire top-20 recommendation set in one shot. This approach
was too brittle for a small local model and frequently produced malformed or
incomplete outputs. The current implementation therefore uses a simpler
two-stage per-item classification strategy: each candidate article is first
placed into a broad group and then refined into one of the final sections
(foundations, core, advanced, adjacent). This does not fully solve the quality
problem, but it produces a more robust prototype behavior than the original
monolithic generation attempt.

Even after this redesign, the limitations of the local organizer remain clear.
Because the current model is small, it sometimes produces semantically weak
classifications or collapses too many articles into similar categories. For
that reason, the present report treats the learning-path layer as an
experimental extension rather than a validated component of the recommendation
system. A larger local model, or a stronger instruction-following model with
better classification consistency, would likely improve this part of the
pipeline substantially.

Two additional solution directions remain active for the next phase:
\begin{itemize}
    \item improving the learning-path organizer with a stronger local model and
    a more reliable sectioning strategy,
    \item extending the revised branch with future experiments on pooling and
    recommendation quality without returning to the earlier flawed split logic.
\end{itemize}
